\section{Method}
\label{sec:method}


We target recurring enterprise workflows in which a task attempt can
modify a shared environment and relevant organization-specific procedures
may be absent from the initial agent harness.
Let $M_r$ denote the persistent memory store at round $r$, $B_r$ a batch
of execution traces, $F_r$ the corresponding environment or human
feedback, and $S$ the ingestion scope governing which traces and feedback
are eligible for consolidation. A consolidation method $C$ updates the
store as
\begin{equation}
    M_{r+1}=C(M_r,B_r,F_r;S).
    \label{eq:consolidation}
\end{equation}
The goal is to capture residual procedural knowledge not already encoded
in the fixed harness. Consolidation transforms many rollout-specific
observations into fewer reusable procedures by merging repeated evidence,
correcting inaccurate knowledge, and retaining distinct lessons when their
conditions or remedies differ. The model parameters, harness, and executor
loop remain fixed; methods differ only in their memory representation and
consolidation operation. After consolidation on training tasks, the store
is frozen and evaluated by the checkpoint and task completion rates the
agent achieves on test tasks.

\subsection{\method{} Consolidation Method}
\label{sec:wasp-method}

\method{} is a memory-consolidation and retrieval method for training-free
continual learning that organizes execution traces and feedback into
persistent, interlinked procedural wiki pages at different granularities and
retrieves them to guide future execution, drawing on the idea of an
LLM-maintained wiki as a personal knowledge base \citep{karpathy2026llmwiki}.
\method{} initializes a wiki filesystem with no initial procedural content and
mandatory control files such as \texttt{index.md}, which serves as a
high-level router over the wiki pages. A \method{} consolidator agent creates,
organizes, revises, or removes pages and their links during consolidation.
Three interface operations are provided:

\begin{itemize}[leftmargin=*,itemsep=2pt,topsep=3pt]
    \item \textbf{Ingestion.} Given a raw agent-session \texttt{.jsonl} file
    and optional asynchronously produced natural-language feedback, the
    \method{} agent spawns a subagent that reviews the session trace to
    identify actionable observations retrospectively, then uses them to add,
    update, or remove pages and links.
    \item \textbf{Query.} During both consolidation and execution, agents begin
    from \texttt{index.md} and combine semantic search with shell navigation to
    locate relevant pages and follow inter-page relationships, such as related
    conditions, actions, and outcomes; the semantic-search index is rebuilt
    after each consolidation phase.
    \item \textbf{Lint.} Prompt-based rules and scripts direct the agent to
    scrutinize generalizability, reusability, ambiguity, excessive page length,
    link integrity, and conflicts, and to determine whether pages should be
    split, merged, or relinked.
\end{itemize}

\paragraph{\method{} variants.}
\emph{Free Form} is the basic \method{} instantiation described above. Each
page is markdown with a minimal metadata header, such as its title, brief
description, and last-modified time; the consolidator decides the remaining
page structure and content. \emph{Schema Form} adds prescriptive,
human-designed organization: pages must be Procedures, Actions, or Lessons,
each with fixed fields and enumerated options. Both \emph{Free Form} and
\emph{Schema Form} read the existing wiki before writing and therefore perform
implicit, agent-managed deduplication. \emph{Schema Form-ED} further adds explicit deduplication
rules: each session is first processed independently into candidate pages,
after which a separate normalization, matching, and merge stage consolidates
candidates before committing them to the persistent store. \emph{Schema Form}
and \emph{Schema Form-ED} thus represent increasing levels of upfront
prescriptive knowledge organization.


\subsection{Controlled Consolidation Design Space}
\label{sec:design-space}

To identify which consolidation choices improve downstream execution and where
human and system effort is most rewarding, we vary four dimensions
systematically while holding the remaining factors fixed.

\paragraph{Consolidation method.}
We compare the three \method{} variants above with four baselines. \emph{No
Memory} executes without assistance from any long-term memory, while \emph{No
Consolidation} stores raw execution traces and retrieves the most similar trace
without abstraction or integration. \emph{Insights} distills each session into
abstract procedural insights \citep{insights} and stores them in a flat
collection, with no cross-session deduplication, verification, or
reconciliation. \emph{SkillX} constructs a knowledge base in the form of agent
skills organized at three levels: planning, functional, and atomic
\citep{skillx}. This hierarchy is prescriptive, analogous to imposing a page
schema: each planning skill forms an independent procedural hierarchy, and
reusable functional or atomic skills are not deduplicated into a shared pool
across those hierarchies.

\paragraph{Feedback type.}
Feedback determines what evidence the consolidator has for deciding whether an
execution was correct and why it succeeded or failed. We construct a ladder
from \emph{Self-Judge}, which uses only the agent's own assessment, to external
\emph{Checkpoint-binary} feedback that localizes success and failure to
individual checkpoints within the task, and \emph{Checkpoint-binary + Human}
feedback that additionally provides natural-language diagnosis. This separates
the value of an external outcome signal from the value of diagnostic
explanation.

\paragraph{Ingestion scope.}
Ingestion scope determines which experiences are eligible for consolidation.
\emph{Success-only} admits evidence from successful executions, whereas
\emph{Both} also admits failed traces and their feedback. Failure traces become
useful when outcome or diagnostic feedback distinguishes a reusable procedural
lesson from incidental unsuccessful behavior.

\paragraph{Refinement rounds.}
Refinement rounds control whether memory is consolidated once or repeatedly
revised as further execution traces become available. R0 serves as the
no-memory baseline, R1 performs the first consolidation, and R2--R4 repeatedly
execute against and update the same persistent store, allowing knowledge to be
corroborated, corrected, or extended over time. Appendix~\ref{app:methods}
specifies the complete controlled matrix and protocols.




